\subsection{拓扑向量空间的向量子空间与商空间；拓扑向量空间的乘积；子空间的拓扑直和}

对拓扑模块 (GT, III, § 6.6) 所说的一切特别适用于拓扑向量空间。若 M 是拓扑向量空间 E 的向量子空间，则由 E 的拓扑在 M 上诱导的拓扑与 M 的向量空间结构相容，并且 M 在 E 中的闭包 \(\overline{M}\) 是 E 的向量子空间。E 的拓扑对 M 的商拓扑与 E/M 的向量空间结构相容。

若 E 是拓扑向量空间，则 \(\{0\}\) 在 E 中的闭包 N（0 的各邻域之交）是 E 的闭向量子空间；商向量空间 E/N 必定是 Hausdorff 的（不论 E 本身是否如此），称为与 E 相伴的 Hausdorff 向量空间。

设 \((\mathrm{E}_{\iota})_{\iota \in I}\) 是同一拓扑除环 K 上的一族拓扑向量空间，并设 E 是诸 \(E_{\iota}\) 的乘积向量空间。诸 \(E_{\iota}\) 的拓扑的乘积拓扑与 E 的向量空间结构相容。在乘积空间 E 中，诸 \(E_{\iota}\) 的直和所构成的子空间 F 处处稠密 (GT, III, § 2.9, prop. 25)。

对 \(\mathbf{R}\) 域或 \(\mathbf{C}\) 域上某些类型的拓扑向量空间，我们 (在 II, 第 29 页) 在一族拓扑向量空间 \((\mathrm{E}_{\mathrm{t}})\) 的直和上定义一个拓扑，它一般不同于由诸 \(\mathrm{E}_{\mathrm{t}}\) 的乘积拓扑所诱导的拓扑。

关于带算子的拓扑群的稳定子群的有限直和所说的一切 (GT, III, §6.2) 都适用于拓扑向量空间，只需把其中的“稳定子群”一律换成“向量子空间”。

注记. --- 给定 Hausdorff 拓扑向量空间 E 的闭向量子空间 M，未必存在 M 在 E 中的（代数）补向量子空间在 E 中闭（即使 E 是赋范空间亦然；参见 IV, 第 55 页, 习题 16 (c)）；更谈不上 M 在 E 中的拓扑补必定存在 (参见 I, 第 26 页, 习题 8)。然而我们将在 § 2 中看到，当 K 是非离散赋值除环时，E 的每个具有有限余维数的闭子空间 M 在 E 中确实有拓扑补 (I, 第 14 页, prop. 3)。

